\documentclass[11pt]{article}

\usepackage{amsmath,amssymb,amsthm,mathtools}
\usepackage[margin=1in]{geometry}
\usepackage{microtype}

\newtheorem{theorem}{Theorem}[section]
\newtheorem{lemma}[theorem]{Lemma}
\newtheorem{corollary}[theorem]{Corollary}
\newtheorem{remark}[theorem]{Remark}

\newcommand{\dd}{\,\mathrm d}

\title{Adjacent Triangle Roots Carrying a Leaf and a Two-Edge Tail:\\
A Second Rooted Supporting-Plane Inequality}
\author{}
\date{}

\begin{document}

\maketitle

\begin{abstract}
Let $H$ be a triangle with one pendant leaf at one triangle vertex and one
internally disjoint two-edge path at an adjacent triangle vertex.  We prove
that every graphon $W$ of edge density $p\geq1/2$ satisfies
\[
 t(H,W)\geq p^4(2p-1)=\Phi_H(p).
\]
The attachments occur at different roots, so the one-root mixed-branch
theorem does not apply.  A two-edge inclusion--exclusion inside the triangle
reduces the density to a scalar function of
$d$ and $a=T_Wd$.  We prove a sharp supporting-plane inequality on the
entire feasible rooted region
$0\leq a\leq d$, $a\geq d-(1-p)$.  Its $a-d^2$ term integrates to zero.
This classifies the formerly open NetworkX Atlas graph $97$.
\end{abstract}


%%%%%%%%%%%%%%%%%%%%%%%%%%%%%%%%%%%%%%%%%%%%%%%%%%%%%%%%%%%%%%%%%%%%%%%%%%%%%%%
\section{The graph, target, and rooted densities}
%%%%%%%%%%%%%%%%%%%%%%%%%%%%%%%%%%%%%%%%%%%%%%%%%%%%%%%%%%%%%%%%%%%%%%%%%%%%%%%

Let $(\Omega,\mu)$ be an arbitrary probability space and let
$W\colon\Omega^2\to[0,1]$ be symmetric and measurable.  Put
\[
 p=\int_{\Omega^2}W(x,y)\dd\mu(x)\dd\mu(y),
 \qquad d=T_W\mathbf1,
 \qquad A=T_Wd,
 \qquad C=T_W(d^2).
\]

Let the triangle vertices be $o,r,s$.  Attach a two-edge path $o-u-v$ at
$o$ and a leaf $w$ at the adjacent triangle vertex $r$, with no other edges.
Integrating the tail and the leaf first gives
\begin{equation}
 t(H,W)=\int_{\Omega^3}
 W(x,y)W(x,z)W(y,z)A(x)d(y)
 \dd\mu(x)\dd\mu(y)\dd\mu(z).
 \label{eq:density}
\end{equation}

The graph has six vertices and six edges.  Once the triangle is coloured,
each new tree vertex has $x-1$ choices after its parent has been coloured.
Therefore
\[
 \chi_H(x)=x(x-1)^4(x-2),
 \qquad \chi(H)=3,
\]
and polynomial continuation gives
\begin{equation}
 \Phi_H(p)=p^4(2p-1).
 \label{eq:target}
\end{equation}
The required interval is $p\in[1/2,1]$.


%%%%%%%%%%%%%%%%%%%%%%%%%%%%%%%%%%%%%%%%%%%%%%%%%%%%%%%%%%%%%%%%%%%%%%%%%%%%%%%
\section{Reduction from two triangle roots to one scalar pair}
%%%%%%%%%%%%%%%%%%%%%%%%%%%%%%%%%%%%%%%%%%%%%%%%%%%%%%%%%%%%%%%%%%%%%%%%%%%%%%%

We reuse two elementary rooted facts, proving them here to keep the argument
self-contained.

\begin{lemma}[Rooted region and star compression]\label{lem:rooted}
For almost every $x$, writing $d=d(x)$ and $a=A(x)$,
\[
 0\leq a\leq d,
 \qquad a\geq d-(1-p),
\]
and
\[
 C(x)\geq
 \begin{cases}
 A(x)^2/d(x),&d(x)>0,\\
 0,&d(x)=0.
 \end{cases}
\]
Moreover,
\[
 \int_\Omega A\dd\mu=\int_\Omega d^2\dd\mu.
\]
\end{lemma}

\begin{proof}
Nonnegativity is immediate, and $A\leq d$ follows from $d(y)\leq1$.
The inequality $W(x,y)d(y)\geq W(x,y)+d(y)-1$ gives
$A(x)\geq d(x)+p-1$.  Cauchy--Schwarz under the finite measure
$W(x,y)\dd\mu(y)$ gives $A(x)^2\leq d(x)C(x)$ when $d(x)>0$.
If $d(x)=0$, then $W(x,y)=0$ for almost every $y$, and $A(x)=C(x)=0$.
Finally, symmetry and Tonelli give
\[
 \int A=\int_{\Omega^2}W(x,y)d(y)\dd\mu^2=\int d^2.
\]
\end{proof}

\begin{lemma}[Adjacent-root reduction]\label{lem:reduction}
Define
\[
 h_p(d,a)=a^2\left(\frac ad+d-1\right)_+
\]
when $d>0$, and $h_p(0,0)=0$.  Then
\[
 t(H,W)\geq\int_\Omega h_p(d(x),A(x))\dd\mu(x).
\]
\end{lemma}

\begin{proof}
For fixed $x$, apply $bc\geq b+c-1$ to $W(x,z),W(y,z)$ in
\eqref{eq:density}.  After multiplication by $W(x,y)d(y)$ and integration
in $y,z$, the three resulting terms are
\[
 d(x)A(x),\qquad C(x),\qquad -A(x).
\]
The inner rooted density in \eqref{eq:density} is also nonnegative.  Hence
\[
 t(H,W)
 \geq\int A(x)\bigl(C(x)-(1-d(x))A(x)\bigr)_+\dd\mu(x).
\]
The positive-part function is nondecreasing.  Lemma~\ref{lem:rooted} gives,
on $\{d>0\}$,
\[
 \bigl(C-(1-d)A\bigr)_+
 \geq A\left(\frac Ad+d-1\right)_+.
\]
On $\{d=0\}$ all relevant quantities vanish.  The assertion follows.
\end{proof}


%%%%%%%%%%%%%%%%%%%%%%%%%%%%%%%%%%%%%%%%%%%%%%%%%%%%%%%%%%%%%%%%%%%%%%%%%%%%%%%
\section{The adjacent-root supporting plane}
%%%%%%%%%%%%%%%%%%%%%%%%%%%%%%%%%%%%%%%%%%%%%%%%%%%%%%%%%%%%%%%%%%%%%%%%%%%%%%%

For $n\geq0$, write
\[
 B_i^n(x)=\binom ni x^i(1-x)^{n-i}
 \qquad(0\leq i\leq n)
\]
for the Bernstein basis on $[0,1]$.  A polynomial expressed in this basis
with nonnegative coefficients is nonnegative on the unit interval; the same
holds for products of Bernstein bases on a box.

\begin{lemma}[Adjacent-root supporting plane]\label{lem:scalar}
Let $p\in[1/2,1]$ and let $d,a\in[0,1]$ satisfy
\[
 0\leq a\leq d,
 \qquad a\geq d-(1-p).
\]
Then
\begin{align}
 h_p(d,a)\geq{}&p^4(3-8p)+2p^3(5p-2)d
 \notag\\
 &\quad+p^2(5p-2)(a-d^2).
 \label{eq:scalar}
\end{align}
\end{lemma}

\begin{proof}
Denote the right side by
\[
 L_p(d,a)=\alpha+\beta d+\gamma(a-d^2),
\]
where
\[
 \alpha=p^4(3-8p),\qquad
 \beta=2p^3(5p-2),\qquad
 \gamma=p^2(5p-2)>0.
\]
If $d=0$, then $a=0$ and $L_p(0,0)=\alpha\leq0$, so assume $d>0$.

Put $t=d(1-d)$.  First suppose $a\leq t$.  Then $h_p(d,a)=0$.
Since $L_p$ is increasing in $a$, it is enough to put $a=t$.  Direct
expansion gives
\[
 L_p(d,t)=-p^2Q_p(d),
\]
where
\[
 Q_p(d)=2(5p-2)d^2+(-10p^2-p+2)d+8p^3-3p^2.
\]
Its leading coefficient is positive and
\[
 \operatorname{disc}_d(Q_p)
 =-(2p-1)^2(5p-2)(11p+2)\leq0.
\]
Thus $Q_p(d)\geq0$, proving \eqref{eq:scalar} in the zero region.

Now suppose $a\geq t$.  Multiplying \eqref{eq:scalar} by $d>0$ reduces it
to
\begin{equation}
 G_p(d,a):=a^3+a^2d^2-a^2d-dL_p(d,a)\geq0.
 \label{eq:G}
\end{equation}
The feasible lower endpoint for $a$ is
\[
 b_p(d)=\max\{d(1-d),d-(1-p)\}.
\]
At its first branch,
\begin{equation}
 G_p(d,d(1-d))=dp^2Q_p(d)\geq0.
 \label{eq:first-boundary}
\end{equation}
The second branch occurs exactly when $d\geq\sqrt{1-p}$.  Put
\[
 C_p(d)=G_p(d,d-(1-p)).
\]
For $p<1$, $C_p$ is a monic quartic and
\[
 \operatorname{disc}_d(C_p)
 =-4p^4(p-1)^4D(p)<0,
\]
where, after $u=p-1/2\geq0$,
\begin{align*}
 D(p)={}&24000u^{16}+123200u^{15}+367520u^{14}+798704u^{13}
 +1276499u^{12}\\
 &+1566514u^{11}+\frac{2961123}{2}u^{10}
 +\frac{2049397}{2}u^9+\frac{8200349}{16}u^8
 +\frac{363635}{2}u^7\\
 &+\frac{1448591}{16}u^6+\frac{43363}{2}u^5
 +\frac{1137901}{256}u^4+\frac{101681}{128}u^3\\
 &+\frac{57211}{512}u^2+\frac{6845}{512}u
 +\frac{3175}{4096}>0.
\end{align*}
Hence $C_p$ has exactly two real roots.  Since
$C_p(0)=(p-1)^3<0$, one is negative and one is positive.

Let $s=\sqrt{1-p}\in(0,1/\sqrt2]$.  Exact factorization gives
\[
 C_p(s)=s(1-s^2)^2R(s),
\]
where
\[
 R(s)=-8s^6-10s^5+11s^4+21s^3-12s^2-9s+5.
\]
To see that $R(s)>0$ without a numerical root estimate, put
$x=4s/3$.  Our interval has $0<x<1$, and the degree-$12$ Bernstein
coefficients of $R(3x/4)$ are
\[
\begin{gathered}
5,\ \frac{71}{16},\ \frac{83}{22},\ \frac{42887}{14080},\
\frac{32447}{14080},\ \frac{72073}{45056},\ \frac{77449}{78848},\\
\frac{22375}{45056},\ \frac{19}{110},\ \frac{1951}{112640},\
\frac{17}{5632},\ \frac{115}{2048},\ \frac{11}{256}.
\end{gathered}
\]
They are all positive.  Therefore $C_p(s)>0$, so the positive real root of
$C_p$ lies below $s$.  It follows that
\begin{equation}
 C_p(d)>0\qquad(d\geq\sqrt{1-p}).
 \label{eq:second-boundary}
\end{equation}
At $p=1$, the same conclusion follows directly from
$C_1(d)=d(d-1)^2(d+5)\geq0$.

It remains only to rule out a dip of the cubic $G_p(d,\cdot)$ below its
nonnegative lower endpoint.  Exact calculation gives
\begin{equation}
 \operatorname{disc}_a(G_p)
 =-d^2p^2(d-p)^2S(p,d),
 \label{eq:cubic-discriminant}
\end{equation}
where $S$ has degree six in $p$ and degree five in $d$.  Put
$x=2p-1$ and $y=d$.  In the product Bernstein basis
$B_i^6(x)B_j^5(y)$, its coefficient matrix $(s_{ij})$ is
\begin{align*}
S(p,d)={}&20d^5p-8d^5-60d^4p+24d^4
 +12d^3p^3-4d^3p^2+60d^3p-24d^3\\
&+1124d^2p^4-924d^2p^3+188d^2p^2-20d^2p+8d^2\\
&-1364dp^5+636dp^4+156dp^3-76dp^2\\
&+1728p^6-1296p^5+243p^4,
\end{align*}
and the coefficient matrix is
\[
\begin{pmatrix}
27/16&97/80&57/80&67/80&83/80&25/16\\
81/16&52/15&517/240&139/60&707/240&9/2\\
1161/80&11783/1200&7837/1200&1537/240&1919/240&2879/240\\
1593/40&1379/50&7801/400&3589/200&1053/50&149/5\\
1053/10&3813/50&8563/150&1523/30&8263/150&713/10\\
270&1032/5&2446/15&431/3&437/3&512/3\\
675&2727/5&2267/5&2017/5&1963/5&423
\end{pmatrix}.
\]
Every entry is at least $57/80$, so $S(p,d)>0$ on the whole box.

If $d\neq p$, \eqref{eq:cubic-discriminant} is negative.  Thus the monic
cubic $G_p(d,\cdot)$ has exactly one real root.  Equations
\eqref{eq:first-boundary}--\eqref{eq:second-boundary} show that this root is
at or below the feasible lower endpoint $b_p(d)$, so $G_p(d,a)\geq0$ for
every feasible $a\geq b_p(d)$.  If $d=p$, direct factorization gives
\[
 G_p(p,a)=(a-p^2)^2(a+3p^2-p)\geq0.
\]
This proves \eqref{eq:G} and the lemma.
\end{proof}

\begin{remark}[Exact polynomial audit]
The large displayed coefficient certificates are independently reconstructed
from the defining polynomials by
\texttt{codes/verify\_triangle\_adjacent\_leaf\_tail.py}.  In particular,
the script does not merely check sampled values; it verifies every identity
over the rational polynomial ring.
\end{remark}


%%%%%%%%%%%%%%%%%%%%%%%%%%%%%%%%%%%%%%%%%%%%%%%%%%%%%%%%%%%%%%%%%%%%%%%%%%%%%%%
\section{The graphon theorem and Atlas consequence}
%%%%%%%%%%%%%%%%%%%%%%%%%%%%%%%%%%%%%%%%%%%%%%%%%%%%%%%%%%%%%%%%%%%%%%%%%%%%%%%

\begin{theorem}[Adjacent leaf and two-edge tail]\label{thm:main}
For every graphon $W$ on an arbitrary probability space, of edge density
$p\in[1/2,1]$,
\[
 \boxed{t(H,W)\geq p^4(2p-1)=\Phi_H(p).}
\]
\end{theorem}

\begin{proof}
By Lemma~\ref{lem:reduction} and then Lemma~\ref{lem:scalar},
\begin{align*}
 t(H,W)
 \geq{}&\int h_p(d,A)\\
 \geq{}&p^4(3-8p)+2p^3(5p-2)\int d
 +p^2(5p-2)\int(A-d^2).
\end{align*}
Lemma~\ref{lem:rooted} says that the last integral vanishes, while
$\int d=p$.  Therefore
\[
 t(H,W)\geq p^4(3-8p)+2p^4(5p-2)=p^4(2p-1),
\]
which is \eqref{eq:target}.
\end{proof}

At $p=1/2$ the target is zero and no division by $2p-1$ occurs.  At $p=1$,
$W=1$ almost everywhere and both sides equal one.  Every balanced complete
$k$-partite graphon, $k\geq2$, gives $d=p$, $A=p^2$, and equality in the
supporting plane, hence realizes equality in the theorem.

\begin{corollary}[Atlas 97]\label{cor:atlas97}
NetworkX Atlas graph $97$, with graph6 code \texttt{EYWO}, is positive.  It
has order six, size six, chromatic number three, and
\[
 \chi_{F_{97}}(x)=x(x-1)^4(x-2).
\]
Every graphon of edge density $p\in[1/2,1]$ satisfies
\[
 t(F_{97},W)\geq p^4(2p-1)=\Phi_{F_{97}}(p).
\]
\end{corollary}

\begin{proof}
Its NetworkX edge list is
\[
 02,\ 12,\ 13,\ 14,\ 24,\ 35.
\]
Vertices $1,2,4$ form a triangle.  The path $1-3-5$ is a two-edge tail at
vertex $1$, while $2-0$ is a leaf at the adjacent triangle vertex $2$.
There are no other edges, so Theorem~\ref{thm:main} applies.
\end{proof}


%%%%%%%%%%%%%%%%%%%%%%%%%%%%%%%%%%%%%%%%%%%%%%%%%%%%%%%%%%%%%%%%%%%%%%%%%%%%%%%
\section{Hidden-assumption audit and exact limitations}
%%%%%%%%%%%%%%%%%%%%%%%%%%%%%%%%%%%%%%%%%%%%%%%%%%%%%%%%%%%%%%%%%%%%%%%%%%%%%%%

All graphon integrands are bounded, nonnegative, and measurable, so every
change in integration order is justified by Tonelli--Fubini.  Pointwise
statements are almost-everywhere statements.  The set $\{d=0\}$ is handled
without division.  The proof uses ordinary homomorphism densities and does
not assume a finite-step graphon, injectivity, regularity, constant degree,
or a positive minimum degree.

For formal verification, the principal new statement is
Lemma~\ref{lem:scalar}.  Bernstein positivity may be formalized by expanding
the two exact coefficient lists and using $B_i^n(x)\geq0$ on $[0,1]$.
The graphon layer needs only the rooted region identities, Cauchy--Schwarz,
the elementary two-edge inequality, monotonicity of positive part, and
integration of the supporting plane.  The graph-specific module awaiting
inspection is \texttt{lean/Taeyoung/Examples/Graph097.lean}.

The theorem requires exactly one leaf and one two-edge path at adjacent
triangle roots.  It does not cover attachments at all three roots, longer
paths, a different number of leaves, or multiple tails.  It also does not
assert a general vertex-gluing theorem: the proof depends on the particular
two-edge inclusion--exclusion and the exact scalar certificate above.

Run the exact algebra audit from the repository root with
\begin{verbatim}
python codes/verify_triangle_adjacent_leaf_tail.py
\end{verbatim}

%%%%%%%%%%%%%%%%%%%%%%%%%%%%%%%%%%%%%%%%%%%%%%%%%%%%%%%%%%%%%%%%%%%%%%%%%%%%%%%
\section{What the Lean formalization actually proves}
%%%%%%%%%%%%%%%%%%%%%%%%%%%%%%%%%%%%%%%%%%%%%%%%%%%%%%%%%%%%%%%%%%%%%%%%%%%%%%%

The machine-checked development is
\texttt{lean/Taeyoung/Methods/AdjTail.lean}, ending at
\texttt{satisfiesLowerBound\_97}.  Theorem~\ref{thm:main} is proved in full, on
all of $p\in[1/2,1]$.  Two things are done differently, and one hypothesis turns
out to be unnecessary.

\paragraph{The departure: no cubic discriminant, no Bernstein certificate.}
This note proves Lemma~\ref{lem:scalar} by treating
$G_p(d,\cdot)=a^3+a^2d^2-a^2d-dL_p$ as a cubic in $a$: it evaluates $G_p$ at
both branches of the feasible lower endpoint --- one of which needs the
degree-$12$ Bernstein expansion of the sextic $R(s)$ in $s=\sqrt{1-p}$ --- and
then rules out a dip by certifying the positivity of
$\operatorname{disc}_a(G_p)$ through an $8\times7$ Bernstein coefficient
matrix.  The formalization certifies none of that.  In the coordinates
\[
 X=a-pd,\qquad Y=d-p
\]
--- both of which vanish at the equality point $(d,a)=(p,p^2)$ \emph{and} at the
degenerate corner $(d,a)=(0,0)$ --- the cleared inequality is the single
\texttt{ring} identity
\begin{equation}
 a^2(a+d^2-d)-d\,L_p(d,a)
 =(a-p^2)^2(a+d^2-d)
 +p^2\bigl(2X^2+(2d+p)XY+(7p-2)dY^2\bigr).
 \label{eq:lean-adj}
\end{equation}
The first summand is nonnegative exactly on the active region $a+d^2-d\geq0$,
where \eqref{eq:G} is what has to be proved.  The bracket is nonnegative on the
whole rooted region $0\leq a\leq d\leq1$, by a split at $d=p/3$.

\emph{Case $d\geq p/3$.}  The form is positive semidefinite outright:
\begin{equation}
 8\bigl(2X^2+(2d+p)XY+(7p-2)dY^2\bigr)
 =\bigl(4X+(2d+p)Y\bigr)^2-q_p(d)\,Y^2,
 \qquad
 q_p(d)=4d^2+(16-52p)d+p^2,
 \label{eq:lean-adj-pd}
\end{equation}
and $q_p\leq0$ on $[p/3,1]$ because $q_p$ is a convex quadratic in $d$ with
\[
 q_p(p/3)=\tfrac{p(48-143p)}9\leq0,
 \qquad
 q_p(1)=p^2-52p+20\leq0
\]
on $p\geq1/2$.  This is the only place the interval endpoint $d\leq1$ is used.

\emph{Case $d\leq p/3$.}  Then $Y<0$, and $0\leq a\leq d$ together with
$p\geq1/2$ give $|X|\leq pd$.  Dropping $2X^2\geq0$ and bounding the cross term
by $-(2d+p)pd(p-d)$ leaves the exact remainder
\[
 (7p-2)dY^2-(2d+p)pd(p-d)=d(p-d)\bigl(6p^2-2p-(9p-2)d\bigr),
\]
whose last factor is at least $p(9p-4)/3\geq0$ once $d\leq p/3$.

So the active case is three \texttt{ring} identities and a handful of sign
checks; no cubic is analysed, no Bernstein basis appears, and no square root of
$1-p$ is ever formed.  The inactive case follows this note exactly: $L_p$ is
increasing in $a$, and at $a=d(1-d)$ it equals $-p^2Q_p(d)$, with the
discriminant bound made explicit as
\[
 8(5p-2)Q_p(d)=\bigl(4(5p-2)d+(2-p-10p^2)\bigr)^2+(2p-1)^2(5p-2)(11p+2).
\]

\paragraph{One hypothesis of Lemma~\ref{lem:rooted} is not needed.}
The feasibility bound $a\geq d-(1-p)$ appears nowhere in the formal proof.  Both
\texttt{plane\_le\_active} and \texttt{plane\_nonpos} hold on the larger region
$0\leq a\leq d\leq1$ alone.  The remaining parts of Lemma~\ref{lem:rooted} are
used and are pre-existing: $A\leq d$ is \texttt{K4Tail.pathOp\_le\_degree},
$A^2\leq d\,C$ is \texttt{Broom.sq\_pathOp\_le}, and $\int A=\int d^2$ is
\texttt{BookTail.integral\_pathOp}.  (The lower bound itself is available as
\texttt{Broom.le\_pathOp} if a later row wants it.)

\paragraph{No positive part, and no division by $d$.}
Lemma~\ref{lem:reduction} is not formalized in the stated form.  Instead of
introducing $h_p$ and reasoning about $(\cdot)_+$ and $A^2/d$, the Lean proof
keeps the inner rooted density
\[
 E(x)=\int_{\Omega^2}W(x,y)W(x,z)W(y,z)\,d(y)\dd\mu(y)\dd\mu(z)
\]
as an object of its own (\texttt{AdjTail.leafTri}), proves the two-edge
inclusion--exclusion in the form
\[
 E\geq C-(1-d)A
 \qquad(\texttt{AdjTail.leafTri\_ge}),
\]
and states the pointwise supporting-plane bound multiplied through by $d$
(\texttt{AdjTail.mul\_plane\_le}).  The set $\{d=0\}$ is then an ordinary case
in \texttt{plane\_le\_rooted} rather than a convention, and the two branches of
$h_p$ become the two branches of a case split on the sign of $a+d^2-d$.  The
identity \eqref{eq:density} is \texttt{AdjTail.homDensity\_adjTail}, obtained
through the generic peeling ladder of \texttt{Methods/Peeling.lean}.

\end{document}
